\chapter{Background and Related Work}
\label{chap:background}

\section{Classical Recommender Systems}
\label{sec:bg-recsys}

Three families dominate e-commerce recommendation: collaborative
filtering (with co-occurrence as its lightweight special case),
content-based filtering, and hybrids of the two. Time-decayed
popularity sits alongside them as the standard method for
trending slots. Each family is covered in detail
elsewhere~\cite{ComprehensiveRecSys2024, LLM4Rec2025}; the formal
definitions used in this thesis are in
Chapter~\ref{chap:foundations}.

\subsection{Co-Occurrence and Collaborative Filtering}
\label{subsec:bg-cooccurrence}

Item-based collaborative filtering rests on a simple pattern:
users who buy product $A$ often also buy product $B$. The core
data structure is a co-purchase matrix that counts the buyers
shared by each pair of items. A similarity score is computed
from it. Jaccard is the most common lightweight choice because
it normalizes by union and is robust to differences in product
popularity~\cite{ComprehensiveRecSys2024}. Item-based variants
win at e-commerce scale because product catalogs have more buyers
per item than items per buyer. This makes the item similarity
matrix denser and more stable. The main cost is precomputation:
production deployments compute the matrix on a schedule, not per
request.

The strengths are interpretability, no model training, and
tolerance to sparse user history. The weaknesses are blindness to
product attributes (two unrelated products that share a buyer
base still link) and the cold-start problem for new products with
no purchase history.

\subsection{Content-Based Filtering}
\label{subsec:bg-content}

Content-based filtering scores similarity using product
attributes instead of purchase signals. A typical setup defines a
feature vector over category, brand, tags, and price proximity,
then computes a weighted distance between the active product and
every other product in the catalog. The main strength is the
zero-purchase cold start: a brand-new product can be recommended
on day one because its attributes are known immediately. The main
weakness is the filter-bubble effect: recommendations reinforce
the active item's profile.

\subsection{Hybrid Recommenders and Trending}
\label{subsec:bg-hybrid}

Burke's taxonomy splits hybrid recommenders into seven patterns;
weighted, switching, and mixed are deployed most
often~\cite{Burke2002Hybrid}. A weighted hybrid combines
normalized scores from multiple algorithms into a single ranking.
A switching hybrid picks which algorithm to run per request based
on context (cold versus warm user, sparse versus rich product). A
mixed hybrid shows results from multiple algorithms side by side
without combining their scores. The prototype uses a switching
hybrid for the personalized ``recommended for you'' slot. Where
scores are combined, the prototype uses Reciprocal Rank Fusion,
which avoids normalizing scores across different algorithms by
combining ranks instead~\cite{LLM4Rec2025}.

Naive popularity over an order log favors products with long-tail
demand and penalizes products that are gaining fast. Exponential
time decay corrects this by giving recent activity more weight,
with a half-life chosen per category: short for fashion and
consumer trends, long for durable
goods~\cite{SessionRecSys2021}. The formal definition is in
Section~\ref{sec:found-similarity}.

\section{Explainable Recommendations}
\label{sec:bg-explainable}

A recommendation with a reason is a different thing from one
without. Classical systems offer pattern-level explanations
(``users who bought X also bought Y'') that are interpretable but
generic. Recent work uses LLMs to generate sentence-level
rationales conditioned on user
context~\cite{LLMEnhancedRecSys2024}. Knowledge-graph-grounded
explanations sit between the two: they draw facts from a curated
graph instead of generating freely, and their rationales contain
fewer hallucinated attributes~\cite{ExplainableRecKG2024}. An
explanation improves trust calibration: users more readily accept
correct recommendations and reject incorrect ones. The thesis
treats explanation as a first-class output of the LLM stage.

\section{Large Language Models in Recommender Systems}
\label{sec:bg-llm-recsys}

Surveys of LLM-augmented recommendation identify three roles for
an LLM in the pipeline: knowledge enhancement (the LLM acts as a
knowledge source for product attributes and user intent),
interaction (the LLM is the conversational interface to a
recommender backbone), and model enhancement (LLM-derived
features feed a downstream ranker as extra
input)~\cite{LLMEnhancedRecSys2024}. The roles differ in cost and
integration depth. Interaction is the cheapest to deploy because
it sits at the surface and needs no retraining. Model enhancement
is the most invasive because it changes the offline training
pipeline. This thesis uses roles 1 and 2: the LLM contributes
knowledge (review themes, comparison narratives, product
explanations) and interaction (the conversational advisor). It
does not retrain a ranker. The data volume of a bachelor-thesis
prototype would not support it, and the architectural lessons of
the thesis are at the integration layer, not the modeling layer.

The standard pattern for combining classical recommenders with
LLMs is a two-stage architecture: a classical candidate generator
returns a small set of products under a tight latency budget, and
an LLM reranks, explains, or summarizes that set in a second
stage~\cite{HSTU2024, ShopifyHSTU}. The pressure is cost: LLM
inference scales with context size, so feeding five candidates is
cheap at consumer scale, while feeding five thousand catalog items
is not. The pattern also separates concerns cleanly. The
classical stage owns relevance and recall over the full catalog;
the LLM stage owns explanation, conflict surfacing, and tradeoff
articulation. The prototype follows this split strictly: the
classical stage runs on the full Convex catalog, and the LLM
stage only ever sees a small candidate set, never the full
catalog.

\section{Review Summarization for Purchase Validation}
\label{sec:bg-review-summarization}

Reviews drive a large share of purchase decisions, yet most
product detail pages show them raw. AI-generated review summaries
compress that reading work into a short structured output:
themes, frequently mentioned features, and (less often) explicit
conflict callouts.

Naive summarization fails in two specific ways. First, it
averages sentiment, hiding the minority opinions that often
matter most for purchase decisions (a 38-versus-9 split
between buyers who love a product and buyers who report failure
after six months is more useful than a 4.2-star average). Second,
when prompted to produce a tidy summary, it tends to fabricate
themes that are not in the source
corpus~\cite{RufusConfidentlyWrong2024}. The second failure has
drawn the most criticism of comparable production review-summary
systems. The thesis treats theme fabrication as a hard
constraint: summaries are generated against the actual review
corpus, themes are validated against the source text before they
are stored, and conflicting opinions are surfaced explicitly
with counts.

\section{AI-Assisted Comparison as Decision Support}
\label{sec:bg-comparison}

Comparison carries the highest cognitive load of any task in a
typical purchase journey. Users compare across mixed attributes
(price, battery life, build material, warranty length, shipping
speed) where the weight of each attribute depends on personal
preferences. LLMs are well suited to laying out these tradeoffs
in natural language. Industry agrees: Shopify's analysis of
agentic shopping flows names comparison and tradeoff articulation
as two of the interactions where LLMs add the most user-perceived
value~\cite{ShopifyAgentic}. Consumer research finds 56\% of
shoppers explicitly want comparison help from AI assistants and
47\% want review summarization~\cite{DigidayAIShopping}.

The main design failure in current AI comparison features is the
``best product'' verdict: the system picks a winner, hides the
tradeoffs, and takes the choice away from the user. UX guidance
for agentic AI explicitly warns against this pattern. It
recommends a tradeoff-first layout where the system surfaces the
structured differences and lets the user weight
them~\cite{SmashingAgenticUX}. The thesis follows the
tradeoff-first pattern: comparisons present structured
differences, never a verdict.

\section{Trust and Human-AI Interaction}
\label{sec:bg-trust}

\subsection{Timing and Consumer Resistance}
\label{subsec:bg-timing}

A CHI 2024 study measured how users react to AI suggestions
delivered at different times. Suggestions offered after a natural
pause were rated significantly more helpful than suggestions
injected during active deliberation~\cite{CHI2024Timing}. Related
work in human-AI personalization makes a stronger claim: reactive
personalization preserves the user's perceived autonomy, and
perceived autonomy is one of the strongest predictors of trust in
AI systems~\cite{ProactiveReactivePersonalization}. Proactive
interventions can be welcome, but only at high-confidence moments
where the system has strong evidence the user is stuck or at
decision risk.

Market data shows consumer trust in AI shopping assistants is
low. 55\% of consumers say they do not trust AI chatbot product
recommendations~\cite{ConsumerDistrust2025}. 42\% see e-commerce
AI assistants as upsell tools, not advisors~\cite{AIUpsell2025}.
71\% of retailers have deployed conversational AI in some form,
but only 16\% of consumers use these assistants
regularly~\cite{ConsumerResistanceAI}. The gap between deployment
and usage points to structural resistance. A trustworthy AI
advisor is not a polish layer on top of a shopping system; the
trust gap is the product question.

\subsection{Behavioral Signals as Intent Indicators}
\label{subsec:bg-behavior-signals}

A long line of information retrieval work shows that lightweight
client-side signals (cursor position, dwell time, scroll depth,
hover trajectory) are reliable proxies for user
attention~\cite{SIGIR2020Cursor}. The signals are cheap to
collect and need no server-side modeling at capture time. Dwell
time deserves a note. In information retrieval, dwell predicts
engagement well: the longer a user reads, the more engaged they
are. In e-commerce, dwell still predicts engagement, but it
predicts intent-to-buy less reliably; users sometimes dwell
because they are confused, comparing across browser tabs, or
distracted~\cite{DwellTimeEngagement2022}. The thesis treats
dwell as an attention proxy, not a buying-intent proxy, and
combines it with review-tab engagement and scroll depth in a
composite readiness score.

A privacy boundary applies. The prototype aggregates raw signals
client-side and persists only summary statistics to the backend.
Raw event streams are never stored or sent off the device.

\section{Case Study: Amazon Rufus}
\label{sec:bg-rufus}

Amazon Rufus is the largest deployed AI shopping assistant by
user count and the most thoroughly studied externally. Rufus is
the reference point for this section: its architecture shows the
practical shape of a production-scale AI shopping system, and its
documented failures inform several design decisions in the
prototype.

\subsection{System Architecture and Capabilities}
\label{subsec:bg-rufus-arch}

Rufus combines three architectural patterns. A multi-model router
picks an inference model per query: simple lookups go to a fast,
cheap model; reasoning-heavy queries go to a slower, more capable
one. The output is two-stream: text is streamed as it generates,
and structured product references are hydrated client-side from
live catalog data, so the user sees the explanation and the live
product card as one coherent
response~\cite{RufusTechnology, RufusScaling}. The system plugs
into Amazon's COSMO knowledge graph, an LLM-built commonsense
graph that enriches products with implicit attributes (a
children's bicycle is also ``a gift for a 7-year-old''),
reportedly improving internal relevance metrics by
60\%~\cite{COSMO}. The integration footprint is broad: Rufus
appears on the product detail page, in search, and in the cart.
It serves over 300 million users with parallel speculative
decoding to keep latency within product
targets~\cite{RufusScaling}.

\subsection{Documented Accuracy and Trust Issues}
\label{subsec:bg-rufus-problems}

External evaluations of Rufus report four classes of failure,
stable across studies. First, recommendation accuracy: independent
measurement places Rufus at 32\% accuracy for picking the right
product for a query, far below the relevance figures Amazon
reports internally~\cite{RufusWrong2024}. Second, factual
hallucination: 28\% of Rufus responses that include a product
price quote the wrong price, an error that is especially
damaging because the user can check it in
seconds~\cite{RufusWrong2024}. Third, self-serving bias: 83\% of
Rufus's recommendations favor Amazon-owned or Amazon-branded
products over third-party alternatives, which reinforces the
perception of the assistant as an upsell
tool~\cite{RufusWrong2024}. Fourth, review summary fabrication:
Rufus has been observed to generate review themes that do not
appear in the source review corpus, the failure that has drawn
the most public criticism of the
four~\cite{RufusConfidentlyWrong2024}.

The four failure modes are not independent. They share a single
root cause: the LLM is allowed to generate factual product data
(prices, attributes, themes) from its parametric memory instead
of reading only from live catalog data and source review text.
This observation motivates the output validation and hydration
strategy in Chapter~\ref{chap:foundations}.

\subsection{Feature Evolution and Lessons}
\label{subsec:bg-rufus-lessons}

Rufus's product timeline is informative beyond the architecture.
The system launched in early 2024 as a reactive assistant: the
user opened a chat surface and the system responded. The
proactive ``Help Me Decide'' surface, where the system offers to
compare or recommend without being asked, was added about ten
months later~\cite{RufusTechnology, RufusScaling}. This order
matches the CHI 2024 timing finding above: reactive came first
because it was easier to launch without breaking user
expectations, and proactive was added piece by piece once the
system had stabilized. Public reporting credits Rufus with
roughly \$12 billion of incremental sales in
2025~\cite{RufusRevenue2025}. The thesis treats this number as a
market-sizing data point, not proof of product-market fit,
because the trust failures above leave the user-experience case
contested.

\section{Positioning of This Thesis}
\label{sec:bg-positioning}

Prior work treats classical recommendation and LLM integration
as separate problems. The questions that decide whether a real
AI shopping advisor is trustworthy (when to speak, how to ground
its answers, how to surface conflicting reviews, how to catch
hallucinations) are scattered across HCI papers, engineering
blogs, and post-mortems of shipped systems.
Table~\ref{tab:bg-positioning} compares this thesis with the
systems from this chapter on the four dimensions where it makes
its design choices: how the AI surface is triggered, where facts
come from, what is checked after generation, and what evidence
backs the configuration.

\begin{table}[ht]
  \centering
  \footnotesize
  \begin{tabular}{p{2.6cm} p{2.2cm} p{2.4cm} p{2.1cm} p{2.1cm}}
    \toprule
    System & Interaction & Grounding & Output checks &
    Config evidence \\
    \midrule
    Classical recommenders \cite{ComprehensiveRecSys2024} &
    none, fixed slots &
    catalog items by construction &
    not needed &
    offline metrics \\
    Shopify generative recs~\cite{ShopifyHSTU} &
    none, ranking only &
    catalog items by construction &
    not needed, no free text &
    A/B tests at scale \\
    KG-grounded explanations~\cite{ExplainableRecKG2024} &
    rationale text per item &
    curated knowledge graph &
    graph lookups &
    offline benchmarks \\
    Amazon Rufus~\cite{RufusTechnology} &
    chat, proactive and reactive &
    RAG over catalog and web &
    none documented &
    internal, not public \\
    This thesis &
    chat, reactive only &
    hydration from live catalog &
    three-check pipeline &
    measured eval sweep \\
    \bottomrule
  \end{tabular}
  \caption{The thesis prototype next to the related systems from
  this chapter.}
  \label{tab:bg-positioning}
\end{table}

The first three rows are safe but limited. Classical
recommenders and Shopify's generative recommender emit only
catalog items, so they cannot hallucinate, but they also cannot
explain, compare, or summarize reviews. Knowledge-graph systems
add rationale text and keep it grounded in a curated graph, at
the cost of building and maintaining that graph. Rufus has the
full conversational surface, but no documented validation layer
between the model and the user, and
Section~\ref{subsec:bg-rufus-problems} lists the failure rates
that follow.

This thesis takes the combination the table leaves open: a
conversational advisor like Rufus, with the by-construction
grounding of the item-only systems, a validation pipeline on
everything the model emits, and a measured evaluation behind
every configuration choice. The contribution is not a new
algorithm. It is an integration architecture and an honest
evaluation against documented production failures, at a scale a
bachelor thesis can defend.
